% !TEX program = xelatex
\documentclass[UTF8,a4paper,zihao=5]{ctexart}
\usepackage[top=2.5cm,bottom=2.5cm,left=2.5cm,right=2.5cm]{geometry}
\usepackage{booktabs,tabularx,longtable}
\usepackage{hyperref}
\hypersetup{hidelinks,pdfauthor={},pdfsubject={全国大学生数学建模竞赛 AI 工具使用详情}}
\setlength{\parindent}{2em}
\linespread{1.25}

\begin{document}
\begin{center}
  {\heiti\zihao{2} AI 工具使用详情}
\end{center}

\section{工具名称和版本}

\begin{longtable}{p{0.18\textwidth}p{0.18\textwidth}p{0.34\textwidth}p{0.16\textwidth}}
  \toprule
  工具名称 & 版本/型号 & 开发机构或公司 & 使用日期 \\
  \midrule
  \textit{待填写} & \textit{待填写} & \textit{待填写} & \textit{待填写} \\
  \bottomrule
\end{longtable}

\section{具体使用目的和环节}

[逐项说明在哪个阶段、为实现什么目的使用 AI，以及影响了论文或代码的哪些部分。不得把核心建模责任转移给 AI。]

\section{关键交互记录}

\subsection{交互一}

\noindent\textbf{提示词：}[保留能够说明任务和约束的重要提示词。]

\noindent\textbf{回复摘要：}[准确概括关键回复；必要时附完整记录。]

\section{采纳和人工修改情况}

[说明哪些建议被采纳、哪些被拒绝，以及如何通过推导、代码、数据或人工审阅进行验证和修改。]

\end{document}
